\section{A common physical middle family for every polynomial}

Let the chain have $L$ sites and local space $\C^q$.
For a region $S$, write $\mathcal C(S)=\{0,\ldots,q-1\}^S$.
Every regional tensor product below lists the complement before the middle.
All intervals are intersected with the chain.

Fix nonnegative integers $a,b,\ell,B,k$ satisfying
\begin{align}
    a+2\ell\le b,\qquad b+2\ell\le L-1,\qquad
    \ell\ge1,\quad q\ge2,\quad k\ge1.
    \label{eq:physical_polynomial_family_hypotheses}
\end{align}
Supply full-chain window operators $W^L_r,W^R_r$, $0\le r<2\ell$,
supported on $[a+r,a+r+2)$ and
$[b+r,b+r+2)$, respectively.
Set $K=[a+2\ell,b+1)$ and supply site-dependent tensors
$A_t(i,j)\in\C^{B\times B}$ for $0\le t<|K|$, together with an arbitrary
closing matrix $C_0\in\C^{B\times B}$. Let $P_M$ be the canonical
full-chain extension of their trace contraction
\begin{align}
    (\sigma,\tau)\longmapsto
    \operatorname{tr}\left(C_0\prod_{t=0}^{|K|-1}A_t(\sigma_t,\tau_t)\right),
    \label{eq:physical_polynomial_supplied_core}
\end{align}
using the increasing enumeration of $K$. Put
\begin{align}
    J          & =[a+\ell,b+\ell+1),                    &
    D          & =(qk)^{18\ell+2\lfloor k/\ell\rfloor}, &
    \mathcal B & =(\max(1,B))^{2k}.
    \label{eq:physical_polynomial_two_bounds}
\end{align}
The count parameter $D$ and the open-boundary bond bound $\mathcal B$
have distinct roles.

These are supporting polynomial-family results for
\cite[Proposition~5.7]{Arad2016RigorousRG}.
The interpolation argument requires only spanning, without independence of
the coefficient operators; see
\cite{gap:aklv13_interpolation_and_word_count}.

\begin{theorem}[The canonical common middle family]
    \label{thm:physical_polynomial_canonical_middle_family}
    \lean{TNLean.MPS.AreaLaw.exists_common_physical_polynomial_mpo_family}
    \leanok
    \uses{def:crossing_frequency_index,thm:global_two_window_polynomial_recovery,
        thm:common_middle_grid_mpo_family,
        thm:common_middle_low_frequency_cardinality,
        def:regional_middle_transferred_family,thm:chain_two_window_family_cardinality}
    Let
    \begin{align}
        \kappa=\{0,\ldots,2\ell-1\}^2\times
        \{0,\ldots,\lfloor k/(2\ell)\rfloor\}^2,
        \qquad m=(k+1)^{4\ell+1}.
        \label{eq:physical_polynomial_crossing_index}
    \end{align}
    For $v=(p,s,i,j)\in\kappa$, put
    $E_v=[a+p+1,b+s+1)$.
    There are nonnegative integers $n_{v,t}$ and matrices $V_{v,t,r}$ in the
    increasing natural coordinates of $E_v$, for $0\le t<m$ and
    $0\le r<n_{v,t}$, which determine the following prescribed family.
    Write $V^{\mathrm{reg}}_{v,t,r}$ for these matrices read on
    $\mathcal C(E_v)$, and set
    \begin{align}
        I_{v,t} & =\{0,\ldots,n_{v,t}-1\}\times
        \mathcal C(E_v\setminus J)^2\times\mathcal C(J\setminus E_v)^2,
                & I                             & =\coprod_{v\in\kappa}\coprod_{t=0}^{m-1} I_{v,t}.
        \label{eq:physical_polynomial_canonical_index}
    \end{align}
    Apply the prescribed regional middle-family transfer to each family
    $V^{\mathrm{reg}}_{v,t}$ from $E_v$ to $J$, and denote the resulting
    matrices by $W_\omega$, $\omega\in I$. Then $|I|\le D^2$, and each
    $W_\omega$, on the paired physical alphabet in increasing site order on
    $J$, has an open-boundary matrix product representation with internal
    bond dimensions at most $\mathcal B$.

    All these choices precede the following data. Supply arbitrary full-chain
    operators $P_L,P_R$ supported on $[0,a+1)$ and
    $[b+2\ell,L)$, and put
    \begin{align}
        H=P_L+P_M+P_R+\sum_{r=0}^{2\ell-1}(W^L_r+W^R_r).
        \label{eq:physical_polynomial_hamiltonian}
    \end{align}
    For every scalar family $c_0,\ldots,c_k\in\C$, there are matrices
    $M_\omega$ on $\mathcal C(J^c)$ such that
    \begin{align}
        \sum_{r=0}^{k}c_rH^r
        =\left(\sum_{\omega\in I}M_\omega\otimes W_\omega\right)^{\mathrm{phys}}.
        \label{eq:physical_polynomial_canonical_decomposition}
    \end{align}
    The superscript denotes the canonical restriction of a full configuration
    to $J^c$ and $J$. No self-adjointness, positivity or gap assumption is made.
    Zero virtual dimension is allowed. This is the general square-derived
    bond bound, with no sharper source bond or running-time assertion.
\end{theorem}
\begin{proof}\leanok
    \uses{thm:global_two_window_polynomial_recovery,thm:common_middle_grid_mpo_family,
        thm:common_middle_low_frequency_cardinality,
        thm:chain_two_window_family_cardinality}
    Choose the common scalar recovery points and weights before the operator
    family and the polynomial coefficients. For every crossing position and
    low-frequency pair, apply the supplied-node middle-family theorem at those
    points. Take the dependent union of its canonical carriers, retaining each
    prescribed regional middle operator. The low-frequency estimates bound the
    number of terms before the final numerical absorption. There are
    $(k+1)^{4\ell+1}\le(k+1)^{4\ell+3}$ scalar nodes, so the square
    cardinality estimate gives $|I|\le D^2$.

    After $P_L,P_R$ and the polynomial coefficients are supplied, scalar
    recovery expresses the polynomial as a sum of the corresponding
    two-frequency coefficients. Choose their exterior matrices in the common
    family and sum them, giving
    (\ref{eq:physical_polynomial_canonical_decomposition}).
\end{proof}

\begin{theorem}[A finite enumeration of the common middle family]
    \label{thm:physical_polynomial_fin_middle_family}
    \lean{TNLean.MPS.AreaLaw.exists_common_physical_polynomial_mpo_family_fin}
    \leanok
    \uses{thm:physical_polynomial_canonical_middle_family}
    Under the same hypotheses and with the same supplied windows, core tensors
    and cap, there are a nonnegative integer $N\le D^2$ and matrices
    $W_0,\ldots,W_{N-1}$ on $\mathcal C(J)$, each with natural open-boundary
    bond bound $\mathcal B$, such that for every supported $P_L,P_R$ and every
    complex family $c_0,\ldots,c_k$ there are exterior matrices $M_r$ with
    \begin{align}
        \sum_{h=0}^{k}c_hH^h
        =\left(\sum_{r=0}^{N-1}M_r\otimes W_r\right)^{\mathrm{phys}}.
        \label{eq:physical_polynomial_fin_decomposition}
    \end{align}
    The number and middle matrices are chosen before both the exterior pair and
    all polynomial coefficients.
\end{theorem}
\begin{proof}\leanok
    \uses{thm:physical_polynomial_canonical_middle_family}
    Enumerate the finite canonical carrier $I$ by
    $\{0,\ldots,|I|-1\}$. Pull both the middle matrices and the later exterior
    matrices back along this bijection. The cardinality, bond bounds and finite
    sum identity are unchanged.
\end{proof}

\begin{theorem}[A common middle family for polynomial evaluation]
    \label{thm:physical_aeval_fin_middle_family}
    \lean{TNLean.MPS.AreaLaw.exists_common_physical_aeval_mpo_family_fin}
    \leanok
    \uses{thm:physical_polynomial_fin_middle_family}
    Under the same hypotheses and fixed data, one may choose $N\le D^2$ and
    middle matrices $W_0,\ldots,W_{N-1}$ with the same natural bond bound
    $\mathcal B$ before every supported exterior pair and every polynomial
    $P\in\C[X]$ of natural degree at most $k$, so that suitable exterior
    matrices satisfy
    \begin{align}
        P(H)=\left(\sum_{r=0}^{N-1}M_r\otimes W_r\right)^{\mathrm{phys}}.
        \label{eq:physical_aeval_fin_decomposition}
    \end{align}
    The zero polynomial is included, with no normalization condition on $P$.
\end{theorem}
\begin{proof}\leanok
    \uses{thm:physical_polynomial_fin_middle_family}
    Apply the preceding theorem with $c_h$ equal to the coefficient of $X^h$
    in $P$. The finite polynomial-evaluation identity gives
    $P(H)=\sum_{h=0}^{k}c_hH^h$. It remains valid in the noncommutative matrix
    algebra because the complex scalar coefficients act centrally.
\end{proof}
